\addcontentsline{toc}{chapter}{Abstract}\vspace{-1cm}
%Border
\begin{tikzpicture}[remember picture, overlay]
  \draw[line width = 4pt] ($(current page.north west) + (0.75in,-0.75in)$) rectangle ($(current page.south east) + (-0.75in,0.75in)$);
\end{tikzpicture}


This work addresses the need for privacy-preserving, verifiable electronic voting in contexts where centralized infrastructure is unreliable or undesirable. Conventional e-voting systems often expose identity–choice linkage through timing correlations, single-vote on-chain writes, or centralized identity handling. These limitations create privacy risks and reduce trust. This report addresses those issues by separating identity tracking from vote recording, using off-chain batching and shuffling with on-chain tallying, and by adopting a local blockchain model suited for edge deployments.

The objectives are to design and implement a voting workflow that preserves anonymity, prevents double voting, and remains auditable. The methodology uses hash-based identity verification, a two-transaction batching model, and a state-driven smart contract that enforces election phases. The system integrates an agentic, voice-guided voting flow to improve accessibility while maintaining strict confirmation logic and neutral guidance.

Simulation and validation are performed on a local blockchain environment to reflect the intended deployment path for modular, edge-device voting terminals. The software stack includes a web-based voting interface, a relayer service for batch submission, and a smart contract that maintains tallies and election state. Test cases include repeated voting attempts, invalid candidates, state transitions, and batch submission thresholds. Results show that identity-to-choice linkage is removed on-chain, double voting is prevented, and tallies remain verifiable through contract state.

Hardware emulation is not performed in this phase; the prototype focuses on software and local-chain validation. Future work includes cryptographic proofs for shuffle correctness, hardened key management, and formal audits to further reduce trust in the relayer and improve deployment readiness.
 

\pagebreak